\section{The Proverbs of Hell}

\begin{quotex}
The road of excess leads to the palace of wisdom.

Truth can never be told so as to be understood, and not be believed.
\flright{\textsc{William Blake}, \emph{Proverbs of Hell}}
\end{quotex}


We owe to \textbf{Rene Guenon} the recognition of a common Tradition. His elucidation of metaphysics and exposition of the underlying symbolism are indispensable to any understanding of Tradition. There is no way forward without a thorough grounding in Guenon's works, keeping in mind that Truth can never be completely and consistently expressed in words, but that gnosis is the only way to grasp it.

Unfortunately, the guenonian path is ultimately a dead end, as few (if any) have decided to follow it. There are those who converted to Islam, as though that were Guenon's real message, while continuing to remain European. Inadvertently, Guenon on the one hand promotes individualism by giving the impression that a man can pick and choose his Tradition and on the other opens the door to a false ecumenism of the New Age or Theosophical sort that he opposed.

This is why the professors will find Guenon so congenial as they misread him as an individualist and a universalist. Far removed, though not so far intellectually, are those neopagans who recreate a pseudo-tradition on the model of a modernist church with a creed to be believed and a moral code to follow.

That is why we reject the notion of multiple traditions. There is only one Tradition, the Hyperborean tradition, that has persisted since the beginning of the world; at least, that is the only tradition that interests us. The method is depth, to penetrate to the unexpressed core of tradition, beyond its varying formulations. This is the only way to understand and not create another system of thought to be believed.

\paragraph{The Road of Excess}


\begin{quotex}
If others had not been foolish, we should be so.
\flright{\textsc{William Blake}, \emph{Proverbs of Hell}}
\end{quotex}

Guenon's exclusive commitment to metaphysics and refusal to address questions of cosmology or contemporary affairs has led to a distortion of our understanding of Tradition. Nevertheless, these questions come up and a metaphysics that somehow floats above the affairs of the word is hardly a complete metaphysics. Anything real has an effect, just as the strong spirit is recognized by its power. Is the call to tradition answered by living the life of an Egyptian peasant or of ``strong warriors? or ``god-inspired priests''?

To his credit, \textbf{Julius Evola} understood this and, by building on Guenon, redefined the nature of tradition. And rather than abandoning the West, he sought to recover and elucidate its deep roots. This task was accomplished neither systematically nor consistently. He dared to consider ideas that were unpopular or even distasteful. The positive result of this is that Evola is forever protected from the claws of the professors who will never be able to sanitize or incorporate him into any academic discipline.

Evola wrote voluminously on many subjects and sometimes ventured into the excessive and even the foolish. Who can accept his belief in telegenesis? One must carefully distinguish between Evola's contributions to our understanding of Tradition and metaphysics on the one hand, and the application of principles to particular historical issues.

For example, while we may adopt as our own the spiritual battle between the forces of \emph{cosmos} and those of \emph{chaos}, although the manifestation of that battle will alter from age to age. After the war, Evola wrote a book criticizing the Fascist state from perspective of the one and only Right, that is, of Tradition. In his essay on cultural Marxism, Evola points out that the founder of Dadaism, \textbf{Tristan Tzara}, was a Jew. What difference does that make when Evola himself had participated in and promoted that movement. Or that \textbf{Schoenberg} and \textbf{Stravinski} were Jews, when death or other metal music is equally lacking in melodic elements. To persist in perpetuating that war will lead only to being blind to the real issues.

Nietzsche wrote somewhere that the errors of great thinkers are of more value than the truths of lesser thinkers. This applies to Evola. If we follow and grasp his thought process, we can gain something, but simply repeating his conclusions or views on contemporary topics of interest will lead to not very much.

\paragraph{The Palace of Wisdom}


\begin{quotex}
No bird soars too high, if he soars with his own wings.
\flright{\textsc{William Blake}, \emph{Proverbs of Hell}}
\end{quotex}


Although Guenon was early in his life involved with Hermetism\footnote{\url{https://gornahoor.net/index.php?s=Hermetism&submit=Search}}, he later criticized it on the grounds that it was only a cosmological doctrine. Besides, he added, it could not be a valid tradition since it served as the basis of paganism, Christianity, and Islam. To the contrary. That argument has value only on the assumption of multiple traditions that arose independently. Yet that is not our position so that the ubiquity of Hermetism demonstrates its value. It is the depth behind western traditions, as we have written many times on Gornahoor.

Hermetism presupposes Traditional metaphysics so there is no contradiction. The Hermetist sees no point in the battle today between paganism and Christianity, properly understood, and regards it as a waste of time. Hermetists of any exoteric path are united in their depth not on superficial allegiances. It can draw on the insights of the Vedanta and Buddhism. It is a system of nobility and power. With its complete cosmology it exceeds Oriental systems because it acknowledges the multiple states of being, states far beyond the yogic realization of the Self.


\flrightit{Posted on 2011-08-23 by Cologero}

\begin{center}* * *\end{center}

\begin{footnotesize}
\begin{sffamily}
\texttt{Perennial on 2011-08-23 at 02:12 said:}

Incredible article! Clears up a lot of the nonsense I see going around, and fits right in with Guenon and the others in their insistence on not mixing traditions. The exoteric traditions must not be confused with Tradition... 

\hfill

\texttt{fa on 2011-09-02 at 09:57 said:}

``Or that Schoenberg and Stravinski were Jews'' hm... = I was never a Jew (Stravinski)

\url{http://andron.sitecity.ru/ltext\_3101202117.phtml?p\_ident=ltext\_3101202117.p\_3101211231}


\hfill

\texttt{Cologero on 2011-09-02 at 10:07 said:}

I suppose, then, that Evola was misinformed. It's a moot point anyway, since after WWII Evola abandoned the anti-semitism that had so motivated him in the 1930s. In his later life, he was even proud of his Dada ``period'', a movement he had at one time condemned as being too ``Jewish''.

\hfill

\texttt{Mr. Fopdoodle on 2017-08-23 at 01:47 said:}

``With its complete cosmology it exceeds Oriental systems because it acknowledges the multiple states of being, states far beyond the yogic realization of the Self.''

Can you elaborate on this? I thought the realization of the unconditioned Self was the end all be all of the spiritual quest--there's nothing higher. If you do this then you're a jivanmukta and there's nothing more to attain, you're finished--you've ``laid down your burden'' and did what had to be done.

But apparently there's more?


\hfill

\end{sffamily}
\end{footnotesize}

